\section{Aligning the lists on a sampled set of layers}
\label{sec:alignment}

The terminal lists $A_I\subseteq M_r$ all have size at most $d$.
For the analytic argument, we need more than disjointness on separated
sets: each label should have a layer common to every set on whose list
it appears. These properties differ even for one label. A label that
occurs on the three sets $\{1,2\}$, $\{1,3\}$, and $\{2,3\}$ has no
common index, although none of these sets is separated from another.

The following abstract lemma supplies the common-index property on a
sampled $s$-element set. Its essential quantifier is that the sampling
law is chosen before the lists, with no dependence on the label universe.
This is what allows the same graph construction to handle every
independent set.

\begin{samepage}
\begin{lemma}[Distributional alignment]
For every pair of integers $s,d\geq 1$ and every real $\eta>0$, there exist
an integer $r\geq s$ and a rational probability distribution $\mu$ on
$\binom{[r]}s$ with the following property.\label{lem:alignment}
Let $M_*$ be any set, and let
$A_I\subseteq M_*$, $|A_I|\leq d$, be specified for every nonempty
$I\subseteq[r]$ with $|I|\leq s$. Assume that
\begin{equation}\label{align:separation}
 A_I\cap A_J=\varnothing
 \qquad\text{whenever }\max I<\min J.
\end{equation}
Then, with probability at least $1-\eta$ over $B\sim\mu$, there exists a map
$a:M_*\to B$ such that
\begin{equation}\label{align:conclusion}
 a(A_I)\subseteq I
 \qquad\text{for every nonempty }I\subseteq B.
\end{equation}
The map $a$ may depend on the list family and on $B$.
\end{lemma}
\end{samepage}

The complete proof is in Section~\ref{sec:alignment-proof}.
It replaces labels by their finite equality patterns, applies minimax
and Ramsey homogenization, and obtains an order-invariant random
pattern. In that limit, separated disjointness forces every label to
have an index common to all of its occurrences. We now apply the
lemma to the lists already extracted from an independent set.

Fix $\eps>0$, and suppose that the layer count $r$ and law $\mu$ in
the graph construction are supplied by Lemma~\ref{lem:alignment}
for $s,d,\eta=\eps$. Section~\ref{sec:parameters} will choose these
data after the alphabet-independent bound $d$ is fixed and before
choosing the Label Cover instance.

\begin{definition}
A pair consisting of a chain and a set $B\in\binom{[r]}s$ is
\emph{aligned} if there is a map $a:M_r\to B$ such that
$a(A_I)\subseteq I$ for every nonempty $I\subseteq B$.
\end{definition}

\begin{lemma}[Aligned pairs occur with high probability]
\label{lem:good-pairs}
Under these choices, for independent uniform chain sampling and
$B\sim\mu$, the probability that the pair is not aligned is at most
$4^r d^2\sigma+\eps$.
\end{lemma}

\begin{proof}
There are at most $4^r$ ordered pairs of subsets of $[r]$.
By Lemma~\ref{lem:separated-lists} and a union bound, the probability
that some separated pair of terminal lists intersects is at most
$4^r d^2\sigma$.
For each chain with no such intersection, Lemma~\ref{lem:alignment}
gives an alignment except for a set of $B$ of $\mu$-probability
at most $\eps$. Adding the two failure probabilities proves the bound.
\end{proof}

The lists were selected before $B$, and include the juntas for every
coefficient choice. An alignment for a fixed aligned pair can therefore
be chosen before the actual coefficient and rotation draws.
The next section uses this fixed alignment to control dependence on
one auxiliary phase for each layer in $B$.
